\documentclass[conference]{IEEEtran}
\IEEEoverridecommandlockouts
% The preceding line is only needed to identify funding in the first
% footnote. If that is unneeded, please comment it out.
%Template version as of 6/27/2024

\usepackage{cite}
\usepackage{amsmath,amssymb,amsfonts}
\usepackage{algorithmic}
\usepackage{graphicx}
\usepackage{textcomp}
\usepackage{amsmath,amssymb}
\usepackage{graphicx}
\usepackage[colorlinks=true, allcolors=blue]{hyperref}
\usepackage[table,xcdraw]{xcolor}
\usepackage{float}
\usepackage{multirow}
\usepackage{caption}
\usepackage{subcaption}
\usepackage{tcolorbox}
\usepackage{booktabs}
\usepackage{subfiles}
\usepackage{import}
\usepackage{adjustbox}
\usepackage{xcolor}
\def\BibTeX{{\rm B\kern-.05em{\sc i\kern-.025em b}\kern-.08em
T\kern-.1667em\lower.7ex\hbox{E}\kern-.125emX}}
\begin{document}

\title{GaitNet: Dynamic, Acyclic Quadruped Gait Generation\\
  %\thanks{Identify applicable funding agency here. If none, delete this.}
}
% \title{Continuous Footstep Optimization for Quadruped Robots via
% Projected Gradient Ascent}

% \author{\IEEEauthorblockN{1\textsuperscript{st} Owen}
%   \IEEEauthorblockA{\textit{dept. name of organization (of Aff.)} \\
%     \textit{name of organization (of Aff.)}\\
%     City, Country \\
%   email address or ORCID}
%   \and
%   \IEEEauthorblockN{2\textsuperscript{nd} Sourav}
%   \IEEEauthorblockA{\textit{dept. name of organization (of Aff.)} \\
%     \textit{name of organization (of Aff.)}\\
%     City, Country \\
%   email address or ORCID}
%   \and
%   \IEEEauthorblockN{3\textsuperscript{rd} Professor Agheli}
%   \IEEEauthorblockA{\textit{dept. name of organization (of Aff.)} \\
%     \textit{name of organization (of Aff.)}\\
%     City, Country \\
%   email address or ORCID}
% }
\author{Owen Sullivan, Sourav Selvaraj, Professor Mahdi Agheli}

\maketitle

\begin{abstract}
  Quadruped locomotion in unstructured environments remains a
  significant challenge. This paper proposes a novel method for
  generating non-gaited, dynamic footstep plans using \textit{GaitNet}. By directly
  exploring the output space through projected gradient ascent,
  GaitNet optimizes footstep coordinates over its own reward
  surface. We formulate the problem with 4 parallel tracks and 4
  steps of projected gradient ascent per foot. Our system is
  evaluated in simulation and a physical Unitree Go1 robot navigating complex
  terrain, demonstrating improvements a state-of-the-art model
  predictive control baseline.
\end{abstract}

\begin{IEEEkeywords}
  Quadruped Locomotion, Footstep Planning
\end{IEEEkeywords}

\section{Introduction}
\label{sec:introduction}

Legged robots enable mobility in
environments that are inaccessible to wheeled and tracked systems.
Their ability to traverse uneven, unstructured, and unpredictable
terrains, such as disaster zones, staircases, dense
vegetation, and rocky landscapes, positions them as key
enablers for future exploration, inspection, and rescue missions.
Unlike wheeled rovers, which are fundamentally limited by the continuous
support surface required for traction, quadrupedal systems can effectively
decouple their path of motion from the continuous topology of the ground,
selecting only the most viable, discrete footholds across otherwise
impassable terrain.

Realizing this potential requires a robot to continuously and
intelligently decide where and when to place its feet, a
process known as \textit{contact planning}. Existing approaches to
this problem broadly fall into three categories: traditional
optimization control, end-to-end learned control, and hybrid
methods that combine the two. Traditional optimization offers strong
guarantees but scales poorly once discrete contact decisions enter
the problem. End-to-end learning solves the mixture of discrete and
continuous decisions implicitly, but sacrifices interpretability and
formal safety. Hybrid methods aim to keep a learned component narrow,
selecting footholds or gait parameters, while a controller based on a
model retains responsibility for stability and tracking. Even
within hybrid designs, most existing planners still commit to a
fixed gait schedule. Contact timing is generated by a periodic phase
schedule or a library of gait primitives, and the learned component
only chooses \textit{where} to place a foot, not whether the
schedule itself should be broken. This constrains a robot's ability
to exploit the full versatility legged locomotion affords on terrain
that a schedule cannot anticipate.

\subsection{Terminology}
\label{subsec:introduction-terminology}

Legged robotics uses several overlapping terms inconsistently across
the literature; we fix their meaning here and use them consistently
throughout. \textit{Footstep planning} refers to the process of
selecting when and where to step. A planner is \textit{greedy} if it
provides only instantaneous actions, with no explicit lookahead over
future footsteps. \textit{Dynamic} locomotion is the opposite of
quasi-static locomotion. Balance is maintained through active control
and dynamic interaction with the environment, rather than by keeping
the center of mass within the support polygon at all times. In this
paper, \textit{dynamic} never means fast or reactive; it refers
specifically to this balance strategy. A \textit{gaited} or
\textit{gait-based} contact sequence is one generated by a periodic
phase schedule or by transitioning between a fixed library of gait
primitives; we call the opposite an \textit{acyclic} gait, a footstep
plan that cannot be represented as such phase cycles or primitive
transitions. \textit{Non-gaited} and \textit{free-gait} are used
synonymously with acyclic elsewhere in the literature; we use
\textit{acyclic} throughout for consistency.

\subsection{Related Work}

\subsubsection{Traditional Control Approaches}
Traditional quadruped locomotion commonly relies on explicit
optimization of footstep positions, contact sequences, and timing via
Model Predictive Control (MPC) \cite{Kim2019highly-dynamic}, which
handles continuous quantities such as joint torques and ground
reaction forces well \cite{Wensing2022Nov} but becomes a
Mixed-Integer Quadratic Program once discrete contact states are
introduced \cite{Geisert2019contact-planning}. Deits and Tedrake
\cite{deits_footstep_2014} formulated footstep placement as an MIQP
over convex terrain regions but restricted motion to sequential
placements without optimizing contact timing; Winkler et al.
\cite{winkler_gait_2018} and Aceituno-Cabezas et al.
\cite{aceituno_simultaneous_2019} extended this to jointly optimize
contact scheduling and timing, at a computational cost that precludes
reactive control in real time. Alternatives based on search, such as
Monte Carlo Tree Search \cite{taouil_non-gaited_2024,
akizhanov_learning_2024}, generate genuinely acyclic sequences, but
high branching factors keep them too slow for correction at high
frequency.

\subsubsection{End-to-End Neural Networks}
End-to-end learning methods map sensor data directly to joint
targets, bypassing explicit dynamics modeling and inherently solving
the mixed-integer contact problem \cite{zhang_learning_2023,
zhang_novel_2024}. This flexibility comes at the cost of formal
stability guarantees, large training data requirements, and
unreliable generalization outside the training distribution.

\subsubsection{Hybrid Methods and Learned Gait Policies}
Hybrid architectures merge learning components with execution based
on a model to balance tractability and stability. Several methods
focus on gait modulation at a high level. Shi et al.
\cite{shi_terrain-aware_2023} modulate trajectory parameters for
efficiency, Xie et al. \cite{xie_glide_2023} train RL policies on
centroidal dynamics given fixed gait heuristics, and blind
proprioceptive approaches \cite{duan_sim--real_2022,
siekmann_blind_2021, lee_learning_2020} push robust traversal without
exteroception. Other systems explicitly select among gaits via Deep
Q-Networks \cite{da_learning_2020}, neural phase integrators
\cite{yang_fast_2021}, or offline gait libraries integrated into
MPC running at high frequency \cite{sun_online_2024}. In every case, the learned
component operates on top of a periodic gait representation and
cannot itself produce contact sequences outside it.

\subsubsection{Learning-Based Footstep Planners}
Closer to our setting, explicit footstep planners driven by data
couple decisions about footholds with an MPC solver. Early architectures such
as FootstepNet \cite{gaspard_footstepnet_2024}, DeepLoco
\cite{peng_deeploco_2017}, and DeepGait \cite{tsounis_deepgait_2020}
laid the groundwork but localized to bipeds or fixed cyclic stepping.
RLOC \cite{gangapurwala_rloc_2022}, Omar et al.
\cite{omar_fast_2022}, and Villarreal et al.
\cite{villarreal_fast_2019} leverage convolutional terrain inputs
combined with fixed or periodic timing. Learned steppability terrain
maps \cite{asselmeier_steppability-informed_2024,
omar_safesteps_2023} inform where to place a foot but not when,
leaving contact timing to an external schedule. Most directly
relevant, ContactNet \cite{bratta_contactnet_2024} uses a
convolutional evaluator to extract non-gaited footstep costs from a
discretized grid, but its selection procedure is restricted to
moving one leg at a time in strict sequence, which caps how far a
system built on it can depart from behavior resembling a gait even
when the underlying cost estimates are acyclic.

\subsection{Motivation and Contributions}
Despite this progress, a clear gap remains. Hybrid planners either
inherit a periodic gait schedule from the control layer below them,
or, like ContactNet, remove the schedule but still commit footsteps
one leg at a time by construction. This motivates our core question:
whether a planner that commits to at most one footstep at a time,
re-deciding faster than a swing completes, can produce dynamic
acyclic gaits that a scheduled or sequential planner cannot.

GaitNet answers this by never representing coordination between legs
explicitly. At each control tick it evaluates a small set of
footstep candidates filtered by terrain per leg, locally refines them
against its own learned value function, and selects at most one
action, or none, for the entire robot. Because a selected swing
outlasts many such ticks, the planner can commit a second leg while
the first is still airborne. Any coordination between legs that
results is an emergent consequence of re-deciding quickly, not
something the architecture was given to reason about, though the
policy may learn to anticipate it implicitly through reinforcement
learning.

Our primary contributions are the following.
\begin{enumerate}
  \item A greedy footstep planner that emits at most one action per
    control tick, from which swing across multiple legs emerges
    without any scheduled phase representation, together with a
    simplified action space that removes the need for the footstep
    evaluation network with two stages used in prior work.
  \item A single-network formulation in which a trained value
    function is reused at deployment time as a local refinement
    objective over a small set of candidates filtered by terrain.
  \item Validation on a Unitree Go1 quadruped against the Legged
    Perceptive Stack, a leading perceptive Nonlinear Model Predictive
    Control (NMPC) baseline, together with ablations of GaitNet
    itself with restricted concurrency that isolate the contribution
    of emergent swing across legs.
\end{enumerate}

This paper builds on our preliminary framework, GaitNet, introduced in
\cite{Sullivan2025}, which established the single-footstep,
single-tick action space and reported promising results in simulation
under a discrete, randomly sampled candidate distribution. Here we
extend that formulation with a continuous candidate search refined by
gradients in place of discrete resampling, and provide the first
hardware validation of the resulting system. Simulation results
attributed to the preliminary framework are reported for context in
\autoref{sec:results}. We have not rerun those experiments under the
system described here, and label such results accordingly throughout.


\section{Methodology}
\label{sec:methodology}
\subsection{System Architecture}
\label{subsec:methodology-system-overview}

This work presents a control framework for quadruped robots capable
of generating dynamic, acyclic gaits in challenging environments
without requiring a predefined periodic schedule. The framework uses
a hierarchical architecture that separates high-level footstep
decision-making from low-level kinematic tracking. At each control
tick, a terrain height scan is thresholded and dilated into a
per-leg validity mask, a small set of candidate footsteps is drawn
within that mask and locally refined against GaitNet's learned value
function, and the single best candidate\textemdash or a no-action
option\textemdash is selected for the entire robot and passed to a
low-level tracking controller for execution via model predictive
control and whole-body feedback. GaitNet's footstep action is
consumed identically regardless of which controller executes it:
in simulation, the tracking controller is the MIT Mini-Cheetah
Controller \cite{di_mini_cheetah_2020}, structured in python by
\cite{zhuang_rl_mpc_locomotion_2025}; on hardware, GaitNet instead
supplies footstep targets directly to the Legged Perceptive Stack's
own NMPC/Whole-Body Control (WBC) pipeline \cite{leggedcontrol}, in
place of that framework's native gait scheduler.

Two constraints are enforced directly in the validity mask, before
GaitNet ever sees a candidate. First, a leg already in its swing
phase is excluded from candidate generation, since it cannot be
redirected mid-swing. Second, a new step is permitted only while at
least \numberstringnum{\corlMinContactLegs} legs remain in ground
contact, which caps concurrent swing at
\numberstringnum{\corlMaxConcurrentSwingLegs} legs; this is a
deliberate feasibility constraint, not an incidental limitation, and
every result in \autoref{sec:results} operates under it. The validity
mask itself is computed from a thresholded height scan, dilated by
\numberstringnum{\corlMaskDilationRounds} rounds of max-pooling to
conservatively account for the physical extent of the robot's foot.

\begin{figure*}[t]
  \centering
  \includegraphics[width=0.7\linewidth]{images/diagrams/control-system.png}
  \caption{System architecture of the proposed control framework. A
    terrain height scan is thresholded and dilated into a per-leg
    validity mask; a small set of candidate footsteps is sampled
    within that mask per leg and locally refined against GaitNet's
    value function; the single best candidate across all legs, or a
    no-action option, is selected and passed to the tracking
  controller.}
  \label{fig:diagram-control-system}
\end{figure*}

\subsection{High-Rate Action Selection}
\label{subsec:methodology-high-rate-selection}

GaitNet is evaluated at \corlControlHz\,Hz, a \corlControlTickMs\,ms
control tick, and never represents multi-leg coordination directly:
at most one footstep is committed per tick, for one leg, and the rest
of the robot is left untouched until its own turn. A selected swing
lasts between \corlSwingDurationMin\,s and \corlSwingDurationMax\,s,
so \numberstringnum{\corlMinTicksPerSwing} to
\numberstringnum{\corlMaxTicksPerSwing} further decision ticks elapse
while that leg is still airborne. GaitNet may therefore commit a
second leg
mid-swing of the first, subject only to the contact-count constraint
above. Multi-leg concurrency is an emergent consequence of
re-deciding faster than a swing completes, not a behavior the
architecture was given machinery to plan; whether the policy learns
to anticipate future ticks and implicitly coordinate legs through
this repeated re-selection is a question of what reinforcement
learning discovers, not something encoded in the action space itself.
This is the central mechanism of the system: the simplicity of a
single-leg, single-tick action space is what makes acyclic, dynamic
gait generation tractable at control-loop rates.

\subsection{GaitNet Formulation}
\label{subsec:methodology-gaitnet-architecture}

GaitNet is a single-network architecture comprised of a robot-state
encoder, a footstep encoder, a shared trunk, and two task-specific
output heads. The first head predicts a scalar value
logit used to rank a candidate; the second predicts the swing
duration to use if that candidate is selected. GaitNet evaluates
candidates conditioned on a one-hot leg encoding (including a
no-action option) and the current tracking state $\mathbf{x_g}$:

\begin{equation}
  \mathbf{x_g} =
  \begin{bmatrix}
    \mathbf p_{b,xy} &
    \mathbf r_{w,z} &
    \mathbf v_b &
    \mathbf \omega_b &
    \mathbf u &
    \mathbf g &
    \mathbf c
  \end{bmatrix}^T
  \label{eq:gaitnet-x}
\end{equation}

Here $\mathbf p_{b,xy} \in \mathbb{R}^{\corlFootPosDim}$ is the $xy$
position of all four end-effectors in the base frame, and
$\mathbf r_{w,z}$ is the height of the robot's center of mass in the
world frame. $\mathbf v_b \in \mathbb{R}^{\corlVelDim}$ and
$\mathbf \omega_b \in \mathbb{R}^{\corlAngVelDim}$ are the robot's
linear and angular velocity in the root frame.
$\mathbf u = [v_x, v_y, \omega_z]$ is the commanded velocity.
$\mathbf g \in \mathbb{R}^{\corlGravityDim}$ is the direction of
gravity in the body frame, and
$\mathbf c \in \{0,1\}^{\corlContactDim}$ is the per-foot contact
state. Together these total \corlStateDim{} dimensions.

The network implements a middle-fusion style feature combinator
\cite{feng2021deep}.
\begin{itemize}
  \item \textbf{Robot State Encoder}: The \corlStateDim-dimensional
    state vector is processed by a
    \numberstringnum{\corlStateEncoderLayers}-layer Multi-Layer
    Perceptron (MLP) with a hidden dimension of
    \corlStateEncoderHidden{} and ReLU activations, producing a
    fixed-dimensional latent state representation.
  \item \textbf{Footstep Encoder}: Each candidate's spatial
    coordinates are embedded by a
    \numberstringnum{\corlFootstepEncoderLayers}-layer feedforward
    network with a hidden size of \corlFootstepEncoderHidden{} and
    ReLU activations. The no-action option is not spatial at all; it
    is routed to a separate, learned embedding vector rather than
    through this encoder.
  \item \textbf{Shared Trunk}: The latent state and the candidate
    embedding are concatenated and passed into a
    \numberstringnum{\corlTrunkLayers}-layer shared trunk, with layer
    normalization applied after each linear mapping.
  \item \textbf{Output Heads}: A value head outputs the unnormalized
    logit $\hat{V}$ used for ranking; a duration head outputs
    $t_{\text{swing}}$, scaled by a sigmoid into
    $(\corlSwingDurationMin, \corlSwingDurationMax)$\,s.
\end{itemize}

In the configuration evaluated in this paper, candidates carry no
learned cost input: the footstep evaluation network used in prior
work \cite{Sullivan2025} to pre-score candidates has been removed
entirely, and GaitNet's own value head is the only source of ranking
information. This choice is motivated directly by an ablation in
\cite{Sullivan2025} in which zeroing that cost input \emph{improved}
survival rate over supplying it (see \autoref{sec:results}),
suggesting the pre-computed heuristic
constrained exploration rather than accelerating it.

\subsection{Candidate Generation and Refinement}
\label{subsec:methodology-gaitnet-gradient-ascent}

For each leg, \numberstringnum{\corlCandidatesPerLeg} candidate
positions are drawn from a Gaussian distribution
($\sigma = \corlCandidateStd$\,m) centered on the leg's nominal
footstep offset, then projected onto the nearest valid cell of that
leg's terrain mask via a precomputed Euclidean distance transform.
Because GaitNet's value head is differentiable with respect to a
candidate's spatial coordinates, each of the
\numberstringnum{\corlCandidatesPerLeg} samples is then refined for
\numberstringnum{\corlGradientAscentIters} iterations of gradient
ascent on $\hat V$ at a learning rate of $\corlGradientAscentLR$, with
the candidate re-projected onto the valid mask after every step:
\begin{equation}
  \mathbf{f}_{k+1} = \Pi_{\mathcal{C}} \left( \mathbf{f}_k + \alpha
  \nabla_{\mathbf{f}} \hat{V}(\mathbf{x}_g, \mathbf{f}_k; \theta) \right)
\end{equation}
where $\Pi_{\mathcal{C}}$ projects onto the nearest valid cell of the
mask $\mathcal{C}$, rather than onto a continuous manifold; the mask
resolution is \corlMaskResolutionCm\,cm, so this step is best
understood as a local search over nearby cells rather than continuous
optimization over the workspace. This refinement is a comparatively
minor part of the pipeline: its purpose is to nudge a candidate toward
better-supported footing within its neighborhood, and the terrain
mask, not the gradient step, is what excludes voids and unsteppable
regions from consideration in the first place.

After refinement, GaitNet evaluates the fixed no-action embedding
alongside the
$\corlCandidatesPerLeg \times \corlNumLegs = \corlTotalCandidatesDeploy$
refined candidates (\numberstringnum{\corlCandidatesPerLeg} per leg,
\numberstringnum{\corlNumLegs} legs) and selects the single
highest-valued option. If no-action wins, or if a leg's mask has no
valid candidates, no footstep is issued that tick.

\subsection{Reinforcement Learning Framework}
\label{subsec:methodology-gaitnet-training}

Because identifying optimal footsteps depends on temporal horizons
containing discrete transitions in the contact state $\mathbf{c}$,
generating stable, labeled ground-truth data for supervised
regression is prohibitively difficult. GaitNet is instead trained
with reinforcement learning, using Proximal Policy Optimization (PPO)
\cite{Schulman_proximal_2017}.

The actor network is GaitNet itself. The critic is a smaller,
separately-sized network that scores each candidate through its own
state encoder, footstep encoder, and trunk, then combines the
per-option scores across all sixteen options with a small combining
MLP to produce a single state-value estimate; it does not share
weights with the actor and omits the duration head. Training uses a
fixed candidate set rather than the refinement procedure above: for
each leg, \numberstringnum{\corlTrainCandidatesPerLeg} candidates are
sampled uniformly at random from the valid terrain mask (plus one
no-action option, \corlTrainTotalCandidates{} candidates total), and
the actor's logits over this set define a categorical
distribution. Gradient-based refinement is used only at deployment
time, applied to this same trained value head; training and
deployment therefore draw candidates from different distributions,
which we treat as an intentional design choice rather than an
oversight, since it lets the value head be learned as a general
scorer rather than overfit to the refinement trajectory.

The duration output for the selected candidate is treated as a
Normal distribution with a fixed standard deviation of
\corlDurationExploreStd\,s, concentrating exploration on the discrete
action choice rather than timing. The total log probability of an
action is the sum of the categorical log probability of the selected
candidate and the normal log probability of its sampled duration.
Reward terms balance velocity tracking against stability and a small
per-step action penalty ($\corlActionPenalty$) that discourages
unnecessary stepping; termination on falling or on stepping into
invalid terrain carries a $\corlTerminationPenalty$ penalty. Full
reward and termination tables, command sampling ranges, and PPO
hyperparameters are given in the appendix. The model is trained in
NVIDIA Isaac Lab \cite{mittal_orbit_2023}, with the physics simulation
and policy running on GPU and the per-robot MPC controllers running in
parallel on CPU, for \corlTrainingIterations{} iterations across
\corlParallelEnvs{} parallel environments, with a terrain-difficulty
curriculum that gradually increases the proportion of unsteppable
terrain as training progresses.


\section{Experimental Results}
\label{sec:results}
% TODO: fill in hardware numbers (Table I, and every XX% / XX.X placeholder
% below) once trial collection against the Legged Perceptive Stack completes.

\subsection{Hardware Comparison}
\label{subsec:results-hardware-comparison}

We evaluate GaitNet on a physical Unitree Go1 against the
\textit{Legged Perceptive Stack} \cite{leggedperceptive}, a
leading framework that is open source, combining NMPC and WBC atop
the OCS2 solver, using rigidly scheduled trajectories optimized in
phase. On hardware, GaitNet supplies footstep targets directly
to this same NMPC/WBC pipeline in place of its native gait scheduler
(\autoref{subsec:methodology-system-overview}), so the comparison
below holds the low-level tracking controller fixed and isolates the
contribution of the planning layer alone.

A physical test track (\autoref{fig:figure-test-environment}) mirrors
the terrain used in simulation. Narrow strips of rigid wooden ground
are interspersed with gaps of varying width, distributed
pseudorandomly. We characterize a track's difficulty as the
percentage of its area that \emph{cannot} be stepped on, ranging from
$\corlDifficultyMin\%$ (flat ground) to $\corlDifficultyMax\%$ (dense
gaps). Because both conditions share a controller, tracking error
would mostly reflect the shared low-level pipeline rather than the
planner. We instead measure overall mission autonomy directly: success
rate (completing the
traversal without the robot's base contacting the ground or drifting
off the track), total steps taken, and traversal time, each over
\corlHwTrialsPerCell{} independent physical trials per method and
difficulty level.

We compare four conditions, isolating the contribution of emergent
concurrency directly:
\begin{enumerate}
  \item \textbf{Legged Perceptive Stack:} the NMPC/WBC baseline
    described above \cite{leggedperceptive}.
  \item \textbf{GaitNet-Single:} GaitNet with the contact-count
    constraint tightened to permit only one leg in swing at a time,
    otherwise identical to the full model.
  \item \textbf{GaitNet-Trot:} GaitNet restricted to only swinging
    diagonally opposing leg pairs concurrently, matching a
    conventional trot's phase structure.
  \item \textbf{GaitNet (Full):} the proposed planner under its
    standard constraint of at most two concurrent swing legs, with no
    restriction on \emph{which} legs those are.
\end{enumerate}
A trial is successful only if the robot completes the full traversal
without falling or placing a foot in a removed terrain section.

\begin{figure}[t]
  \centering
  \includegraphics[width=\linewidth]{images/figures/test-environment.png}
  \caption{The Unitree Go1 platform navigating the physical testing
    track used for hardware evaluation. Terrain difficulty is the
    percentage of the track area that cannot be stepped on.}
  \label{fig:figure-test-environment}
\end{figure}

\begin{table}[t]
  \centering
  \caption{Hardware Metrics on the Unitree Go1, \corlHwTrialsPerCell{}
    Trials per Method}
  \small
  \begin{tabular}{lccc}
    \toprule
    \textbf{Method} & \textbf{Success} & \textbf{Time [s]} & \textbf{Steps} \\
    \midrule
    Legged Perceptive Stack \cite{leggedperceptive} & \corlHwLPSSuccess\% & \corlHwLPSTime{} & \corlHwLPSSteps{} \\
    GaitNet-Single & \corlHwSingleSuccess\% & \corlHwSingleTime{} & \corlHwSingleSteps{} \\
    GaitNet-Trot & \corlHwTrotSuccess\% & \corlHwTrotTime{} & \corlHwTrotSteps{} \\
    \textbf{GaitNet (Full)} & \textbf{\corlHwFullSuccess\%} & \textbf{\corlHwFullTime{}} & \textbf{\corlHwFullSteps{}} \\
    \bottomrule
  \end{tabular}
  \label{tab:hardware-metrics}
\end{table}

\autoref{tab:hardware-metrics} and \autoref{fig:data-hardware-survival}
summarize hardware results across the full difficulty range.
GaitNet (Full) achieves the highest success rate of the four methods,
with the gap widening substantially at high difficulty
($\corlHighDifficultyRangeLow$--$\corlDifficultyMax\%$).
The Legged Perceptive Stack's fixed phase schedule, and the
GaitNet-Trot and GaitNet-Single ablations, all increasingly commit to
footholds inside a scheduled or single-leg sequence that cannot be
revised once the terrain beneath it turns out to be invalid, whereas
GaitNet (Full) reevaluates every leg independently at every 40\,ms
tick and can commit a second leg mid-swing to recover from a
foothold the first leg is already unable to use productively. This
matches the ordering predicted by
\autoref{subsec:methodology-high-rate-selection}: GaitNet-Single, the
most constrained ablation, performs closest to the scheduled baseline,
GaitNet-Trot recovers some concurrency and some performance, and
GaitNet (Full) is uniformly best.

\begin{figure}[t]
  \centering
  \includegraphics[width=\linewidth]{images/data/experiments/Gaitnet
  - survival-curr - 69.4 - individual.png}
  \caption{Hardware success rate versus terrain difficulty for all
    four evaluated methods. \textit{Placeholder: simulation figure
    shown pending hardware data collection.}}
  \label{fig:data-hardware-survival}
\end{figure}

\subsection{Simulation Results (Preliminary Framework)}
\label{subsec:results-simulation}

Before developing the system described in
\autoref{sec:methodology}, our preliminary framework \cite{Sullivan2025}
evaluated cost-ablated GaitNet against an independent single-leg
baseline planner in NVIDIA Isaac Lab, across a grid of terrain
difficulties ($\corlDifficultyMin\%$ to $\corlDifficultyMax\%$, in
$\corlDifficultyStep\%$ increments) and commanded forward velocities
($\corlVelocityMin$ to $\corlVelocityMax$\,m/s), over
\corlSimTrialsPerCell{} trials of \corlSimEpisodeDuration\,s each per
cell. These results were collected under the discrete,
randomly sampled candidate distribution described in
\autoref{subsec:methodology-gaitnet-training}
(\corlTrainCandidatesPerLeg{} candidates per leg, cost-ablated) rather
than the continuous candidates refined by gradients used at
deployment in this paper
(\autoref{subsec:methodology-gaitnet-gradient-ascent}). We report
them here for context. We have not rerun this grid under the
refinement procedure introduced in this paper, so they should not be
read as validating the system described here.

In that configuration, GaitNet (Full) achieved a
$\corlResCostAblatedSurvivalRate\%$ mean survival rate, compared to
$\corlResBaselineSurvivalRate\%$ for the independent single-leg
baseline, a footstep selector using the lowest cost, operating one
leg at a time
with a fixed $\corlBaselineSwingDurationMs$\,ms swing duration,
structurally similar to the GaitNet-Single ablation introduced above
but without a learned value function. GaitNet maintained strong
performance for commanded velocities below
$\corlStrongPerfVelocityThreshold$\,m/s, while the single-leg
baseline's performance degraded once commanded velocity exceeded
$\corlVelocityMin$\,m/s and degraded further as terrain difficulty
increased. The performance gap narrowed only in the easiest condition
($\corlVelocityMin$\,m/s, flat ground), consistent with concurrent,
high-rate footstep selection becoming more valuable as task difficulty
increases.

\subsection{Emergent Concurrency Analysis (Preliminary Framework)}
\label{subsec:results-concurrency}

The central claim of this paper is that coordination between legs
emerges from high-rate selection of a single action rather than being
planned. Our preliminary framework \cite{Sullivan2025} first
demonstrated this mechanism in simulation, under the same discrete
candidate distribution described in
\autoref{subsec:results-simulation}. \autoref{fig:data-swing-schedule}
reproduces a representative swing schedule from that study. Below a
$\corlSwingScheduleTransitionMark$\,s mark, at a low commanded
velocity, the robot moves one foot at a time with no overlap between
swings. After the commanded velocity increases, the robot begins
committing a second leg while the first is still airborne, without a
strict alternating pattern of any kind. This transition is not
triggered by any explicit mode switch in GaitNet. The same value
function and the same per-tick decision rule govern both regimes, and
the increased concurrency emerges purely from more candidates crossing
the no-action threshold more often as commanded velocity rises.
Because this decision rule is unchanged by the continuous refinement
procedure introduced in this paper (\autoref{sec:methodology}), we
expect the same emergent behavior to hold at deployment. The hardware
comparison in \autoref{subsec:results-hardware-comparison} tests this
directly.

\begin{figure*}[t]
  \centering
  \includegraphics[width=0.85\textwidth]{images/data/swing-schedule.png}
  \caption{Example swing schedule. Each row represents a leg, with
    color indicating swing phase. Below
    \corlSwingScheduleTransitionMark\,s the robot is commanded at low
    velocity and steps one foot at a time; above
    \corlSwingScheduleTransitionMark\,s, at higher commanded velocity,
    swings overlap without a fixed alternating pattern.}
  \label{fig:data-swing-schedule}
\end{figure*}

In this same study, swing durations vary rather than following a
fixed schedule. \autoref{fig:data-swing-duration-histogram} shows a
peak between
$\corlSwingHistPeakMin$\,s and $\corlSwingHistPeakMax$\,s for typical
forward motion, with a secondary mode down to
$\corlSwingHistSecondaryMode$\,s used for situational, faster steps
such as recovering a marginal foothold. The mean swing duration
across training settles to $\corlResMeanSwingDuration$\,s with a
standard deviation of $\corlResStdSwingDuration$\,s, indicating the
duration head is in active use rather than collapsing to a constant.

\begin{figure}[t]
  \centering
  \includegraphics[width=\linewidth]{images/data/swing_duration_histogram.png}
  \caption{Histogram of swing durations selected by GaitNet during
  evaluation.}
  \label{fig:data-swing-duration-histogram}
\end{figure}

\subsection{Swing Duration Ablation (Preliminary Framework)}
\label{subsec:results-duration-ablation}

To assess whether this dynamic duration output is important, our
preliminary framework \cite{Sullivan2025} compared standard GaitNet
against a Duration-Ablated-GaitNet trained identically but constrained
to a fixed swing duration of $\corlResDurationAblationFixedDuration$\,s,
the approximate mean of the standard model's learned distribution. The
result was a null finding. The duration-ablated model achieved a
$\corlResDurationAblatedSurvivalRate\%$ survival rate, statistically
indistinguishable from the $\corlResGaitNetSurvivalRate\%$ achieved by
the corresponding standard GaitNet variant under matched conditions. We
do not read this as evidence that dynamic duration is unimportant in
general. More plausibly, the tested terrain and velocity ranges were
not demanding enough to expose a difference, and the reward function
does not explicitly incentivize using the full duration range beyond
what is needed to avoid termination. We report this honestly as a
limitation of that evaluation rather than removing the duration head
from the architecture used in this paper, since its cost is negligible
and it may prove useful at higher commanded velocities or on more
aggressive terrain than tested there.


\section{Conclusion}
\label{sec:conclusion}
This paper presented GaitNet, a footstep planner that generates
dynamic, acyclic quadruped gaits without representing multi-leg
coordination anywhere in its architecture. At each \corlControlTickMs\,ms
control tick, GaitNet commits at most one footstep, drawn from a small
set of terrain-filtered candidates and locally refined against its own
learned value function; because a swing outlasts several such ticks,
multi-leg concurrency emerges from high-rate re-decision rather than
from an explicit schedule. On physical hardware, GaitNet (Full)
outperforms both the Legged Perceptive Stack, a state-of-the-art
perceptive NMPC baseline, and concurrency-restricted ablations of
itself (GaitNet-Single, GaitNet-Trot), achieving
\textbf{\corlHwFullSuccess\%} versus \textbf{\corlHwLPSSuccess\%} mean
success rate across matched terrain difficulties, with the gap
widening at the highest difficulties tested. Simulation
results and a swing-schedule analysis corroborate the mechanism
directly: the same policy moves one foot at a time under easy
commands and produces overlapping, non-alternating swings under
harder ones, with no change to the decision rule that produces either
behavior.

\subsection{Limitations and Future Work}

GaitNet's concurrency is capped at
\numberstringnum{\corlMaxConcurrentSwingLegs} simultaneous swing legs
by an explicit feasibility constraint requiring at least
\numberstringnum{\corlMinContactLegs} legs in ground contact; this was
a deliberate design choice for stability margin, but its effect on
performance at the extremes of terrain difficulty is untested and
worth ablating directly. The policy is also trained on a candidate
distribution, \numberstringnum{\corlTrainCandidatesPerLeg} randomly
sampled positions per leg, that differs from the distribution it is
deployed against, \numberstringnum{\corlCandidatesPerLeg} candidates
locally refined by gradient ascent; we treat this as intentional but
have not yet characterized how much the refinement step contributes
beyond what random sampling alone would achieve. Both simulation and hardware deployments depend on an
external height scan to build the terrain validity mask, and the two
platforms currently source that scan differently (a raycast height
field in simulation, the Legged Perceptive Stack's heightmap on
hardware), which is a source of sim-to-real gap we have not yet
quantified.

Future work will extend the state representation with direct
exteroceptive input, such as point clouds or height maps processed
by convolutional layers, to reduce dependence on hand-tuned terrain
thresholding, and will expand the action parameterization to include
whole-body posture and center-of-mass placement, enabling the same
high-rate re-decision principle to extend to steep inclines and
stair climbing without hard-coded elevation heuristics.


\bibliographystyle{ieeetr}
\bibliography{refs}

\end{document}
